\Needspace{18\baselineskip}
\section{Concrete focuses and the boundary argument of Lemma 39(1)}
\label{sec:concrete-focus}

This section returns to the proof order of Appendix~F. We formalize the
boundary argument in Proof~F.4: replacing focus-free fragments preserves the
context's focus mode and whether its source is an input. The result holds for
arbitrarily deep sites and finite sequences of fillings. It does not use
output-closure, rootedness, determinism, or transitions. It supplies the
initial/final agreement needed in Lemma~39(1), for the present substitution
representation. Connecting it to the paper's full specimen definition is still
an obligation; no concrete \texttt{SpecimenSystem} is constructed here.

\subsection{Global focus freedom and a unique shallow focus}
\begin{definition}[Focus labels and focus freedom]\label{term:focus-free}
For a mode $m$, the symbol $\mathsf{focusLabel}(d,m)$ is the primitive focus
symbol in the alphabet of depth $d$. At depth zero it uses the universe-lift
wrapper; at a successor depth it uses the primitive summand.
A label is \emph{focus-free} if it is a primitive other than every focus symbol,
or a box whose body has only focus-free labels. A graph is focus-free when all
its edge labels satisfy this recursively defined condition. Thus the condition
excludes focuses inside boxes at every depth, not merely at the outer layer.
\end{definition}
\lean{Graph/Focused}{13}{15}{Linear.focusLabel}
\lean{Graph/Focused}{18}{21}{Linear.LabelFocusFree}
\lean{Graph/Focused}{24}{25}{Linear.FocusFree}

\begin{definition}[Focussed hypernet]\label{term:focussed-hypernet}
A focussed hypernet consists of a hypernet $G$, a distinguished shallow edge
$f$, its mode $m$, and the following certificates. The label of $f$ is
$\mathsf{focusLabel}(d,m)$; every other edge has a focus-free label; the source
and target lists of $f$ are respectively $[s]$ and $[t]$; $s\ne t$; and it is
not the case that $s$ is an input and $t$ an output.
The label of $f$ is primitive, so has no box children. Consequently these
certificates enforce exactly one focus globally, with the shallowness,
self-acyclicity, and non-exposure required in \S6.2. The explicit singleton
incidence certificates record the endpoints without choosing positional casts.
\end{definition}
\lean{Graph/Focused}{29}{39}{Linear.Focussed}

\begin{lemma}[Focus-free addressed substitution]
If $G$ and $H$ are focus-free, so is $\mathsf{plugAt}(s,H)$ for every site $s$ in $G$
with the interface of $H$.
\end{lemma}
\begin{proof}
Induct on the site. At a shallow site, every surviving edge retains its old
label, and every inserted edge has a label of $H$, so both kinds are focus-free.
Inside a box, focus freedom of the host implies focus freedom of that box's
body. Apply the induction hypothesis to the body. Replacing its label preserves
focus freedom of that edge; every other outer label remains unchanged.
\end{proof}
\lean{Graph/Focused}{42}{63}{Linear.focusFree_plugAt}

\begin{definition}[Boundary initiality and finality]\label{term:boundary-finality}
For a focussed hypernet, define $\mathsf{Initial}$ to mean that the mode is
search and the focus source is an input. Define $\mathsf{Final}$ by replacing
search with backtrack. These predicates are defined on any interface for use
in the preservation lemmas. On states of type $*\to\epsilon$, there is exactly
one input, giving the boundary characterizations used in \S6.4 and Proof~F.4.
Finality is a syntactic predicate here; absence of outgoing transitions from
final states has not yet been proved, because transitions on this explicit focus-edge state representation are not yet
defined. Section~\ref{sec:search} proves the search-only boundary laws in coordinates.
\end{definition}
\lean{Graph/Focused}{69}{70}{Linear.Focussed.Initial}
\lean{Graph/Focused}{73}{74}{Linear.Focussed.Final}

\subsection{Preservation through one filling}
\begin{lemma}[Shallow filling preserves the focussed structure]
Let $G$ be focussed at $f$ and let $e\ne f$. Replacing $e$ by a focus-free $H$
produces a focussed hypernet with the same mode.
\end{lemma}
\begin{proof}
The new focus is the retained edge $\mathsf{inl}(f)$ with unchanged label.
Every other retained edge inherits focus freedom from $G$; every inserted edge
inherits it from $H$. Its endpoints are the host images of $s,t$, and its
incidence lists are their singleton lists by functoriality of list mapping.
The host map is injective, so these endpoints remain distinct. It maps internal
ports to internal ports and fixes input and output constructors. Therefore
exposure cannot be introduced.
\end{proof}
\lean{Graph/Focused}{77}{107}{Linear.Focussed.plug}

\begin{lemma}[Shallow boundary preservation and agreement]
A shallow filling as above is initial exactly when its host is initial, and
final exactly when its host is final. Any two focus-free fillings of the same
selected edge therefore agree on initiality and finality.
\end{lemma}
\begin{proof}
The focus mode is unchanged. The host image of its source is an input if and
only if the original source is an input, by case analysis on the three port
constructors. Substitute these facts into the two boundary predicates. For two
fillings, compose each equivalence with the inverse equivalence for the other
filling. This is exactly the argument of Proof~F.4.
\end{proof}
\lean{Graph/Focused}{111}{116}{Linear.Focussed.plug_initial_iff}
\lean{Graph/Focused}{118}{123}{Linear.Focussed.plug_final_iff}
\lean{Graph/Focused}{127}{135}{Linear.Focussed.specimen_boundary_agreement}

\begin{lemma}[Replacing a box body preserves an outer focus]
Replacing a box body by a focus-free body of the same interface preserves a
surrounding focussed hypernet.
\end{lemma}
\begin{proof}
The box edge cannot be the focus edge: their labels belong to distinct summands
of the alphabet. Keep the focus, its mode, and its two endpoints. All incidence
lists are unchanged by box-body replacement, so the self-acyclicity and
non-exposure certificates remain valid. The changed label is focus-free by the
hypothesis on the new body, and all other non-focus labels are unchanged.
\end{proof}
\lean{Graph/Focused}{138}{164}{Linear.Focussed.replaceBox}

\begin{definition}[A site avoiding the focus]\label{term:focus-avoiding-site}
A shallow site avoids the focus when its selected edge differs from the focus
edge. A site strictly inside a box automatically avoids the focus, since the
focus is shallow and primitive. This condition does not assert that the
selected edge is a hole: the preservation results intentionally apply to a
larger class of replacements than paper specimens require.
\end{definition}
\lean{Graph/Focused}{168}{172}{Linear.Focussed.Avoids}

\begin{lemma}[Deep focussed filling and boundary agreement]
A focus-free filling of any site avoiding the focus preserves the focussed
structure and both boundary predicates. Consequently two such fillings of one
site agree on initiality and finality.
\end{lemma}
\begin{proof}
For a shallow site use the shallow filling construction. For a site inside a
box, the box is distinct from the focus, hence its old body is focus-free.
Focus-free addressed substitution proves the modified body focus-free. Apply
the box-body preservation construction. In this second case the outer mode and
endpoints remain literally the same, so both boundary equivalences are
reflexive. Finally compose the equivalences for two fillings through their
common host.
\end{proof}
\lean{Graph/Focused}{175}{190}{Linear.Focussed.plugAt}
\lean{Graph/Focused}{193}{200}{Linear.Focussed.plugAt_boundary_iff}
\lean{Graph/Focused}{204}{211}{Linear.Focussed.deep_specimen_boundary_agreement}

\subsection{Finite fillings and the state space}
\begin{definition}[Finite filling sequence]\label{term:focus-fillings}
A filling sequence starts with a focussed graph and consists of finitely many
addressed substitutions of focus-free graphs at sites avoiding the current
focus. The constructors are the empty sequence and extension by one filling.
Sites may occur at any box depth. Each intermediate focus is the certified
focus transported by the preceding construction. In particular, successive
fillings of a context's holes are included whenever their sites are supplied.
This representation does not yet define a simultaneous tuple-substitution API
or identify all such sequences modulo permutation.
\end{definition}
\lean{Graph/Focused}{215}{221}{Linear.Focussed.Fillings}

\begin{theorem}[Boundary agreement for finite fillings]
Every finite filling sequence preserves initiality and finality. If two
sequences start from the same focussed context, their endpoints agree on both
predicates.
\end{theorem}
\begin{proof}
Induct on the sequence. The empty case is reflexivity. In the extension case,
compose the one-step boundary equivalence with the induction hypothesis. For
two sequences, compose the resulting equivalences through their common initial
context. No template relation between replacements is necessary. Thus once
concrete specimens supply these filling witnesses, the initial/final conclusion
of Lemma~39(1) follows without an extra assumed locality law.
\end{proof}
\lean{Graph/Focused}{224}{231}{Linear.Focussed.fillings_boundary_iff}
\lean{Graph/Focused}{235}{241}{Linear.Focussed.fillings_boundary_agreement}

\begin{definition}[Concrete state]\label{term:concrete-state}
A state packages a hypernet of type $*\to\epsilon$, a focussed structure on that
hypernet, and a proof that it contains no holes at any depth. This is the state
space of \S6.4 in our graph representation. It is not yet equipped with a
\texttt{Machine}: the coordinate search relation of Section~\ref{sec:search} still needs a realization
bridge, and substitution and compute transitions remain open.
\end{definition}
\lean{Graph/Focused}{247}{250}{Linear.State}

\subsection{Nonvacuous checks}
The regression graph has an input, an internal vertex, and an output, with a
focus edge from the input to the internal vertex followed by a unary hole.
It is obtained by relabelling the existing two-edge chain, preserving all
signatures. Its focus is certified directly: the only other edge is the hole,
the endpoints differ, and its target is internal, so the focus is not exposed.
\lean{Graph/FocusedExamples}{9}{13}{Linear.Examples.focusChain}
\lean{Graph/FocusedExamples}{15}{34}{Linear.Examples.focusChain_focussed}
The unary variable corolla is focus-free because a variable label cannot equal
a focus label. Filling the second edge of the search chain by this corolla
remains initial by shallow boundary preservation. The backtracking chain is
final, whereas the search chain is not, since the two mode constructors differ.
\lean{Graph/FocusedExamples}{36}{38}{Linear.Examples.variableFragment_focusFree}
\lean{Graph/FocusedExamples}{41}{46}{Linear.Examples.search_filling_initial}
\lean{Graph/FocusedExamples}{49}{51}{Linear.Examples.boundary_modes_distinct}
Finally, boxing the search chain gives a graph whose outer edge is a box but
whose body contains a focus. If it were focus-free, applying that predicate to
the outer box edge and then the first body edge would say that the search focus
label differs from itself, a contradiction. This checks that the definition
really inspects nested boxes.
\lean{Graph/FocusedExamples}{55}{60}{Linear.Examples.boxed_focus_not_free}
